\documentclass[11pt]{article}

\usepackage[margin=1in]{geometry}
\usepackage{amsmath,amssymb}
\usepackage{array}
\usepackage{enumitem}

\title{Lecture 2 Notes: \\ A Potpourri of Questions and Attempted Answers}
\author{Philosophy of Mathematics}
\date{}

\begin{document}

\maketitle

\section*{Central Theme}

A philosophy of mathematics must explain how mathematics can be necessary,
knowable, objective, meaningful, and applicable.

\[
\boxed{
\text{Philosophy of mathematics explains the place of mathematics in inquiry.}
}
\]

The guiding question for the lecture:

\[
\boxed{
\text{What must a philosophy of mathematics account for?}
}
\]

This chapter surveys the main problems rather than defending one final answer.

\section*{1. Mathematics as an Epistemological Problem}

Mathematics is involved in many of our best attempts to gain knowledge. But mathematical
knowledge seems different from ordinary scientific knowledge.

Scientific claims appear contingent and revisable:

\[
\text{There are nine planets.}
\]

\[
\text{Gravity obeys an inverse-square law.}
\]

These claims might have been false. Scientific history also contains revolutions in which
central claims were rejected.

Mathematical claims seem different:

\[
7+5=12
\]

\[
\text{There are infinitely many primes.}
\]

\[
\text{The induction principle holds for the natural numbers.}
\]

These do not seem merely probable. They seem necessary.

\subsection*{Teaching Point}

Mathematics seems to give us knowledge that is more secure than ordinary empirical
knowledge. A philosophy of mathematics must explain this appearance.

\subsection*{Whiteboard}

\[
\begin{array}{c|c}
\textbf{Science} & \textbf{Mathematics} \\
\hline
\text{contingent} & \text{necessary?} \\
\text{empirical} & \text{a priori?} \\
\text{probabilistic} & \text{proof-based} \\
\text{revisable} & \text{certain?}
\end{array}
\]

\section*{2. Necessity and A Priori Knowledge}

The chapter introduces two central notions.

\[
\textbf{Necessity:}
\quad
\text{things could not have been otherwise.}
\]

\[
\textbf{A priori knowledge:}
\quad
\text{knowledge independent of experience.}
\]

Mathematics has traditionally been treated as a paradigm of both.

\[
\text{Mathematics} = \text{necessary} + \text{a priori}
\]

But both notions require explanation. It is not enough simply to say that mathematics is
necessary and a priori. We need an account of what these terms mean and why they apply
to mathematics.

\subsection*{Teaching Point}

The traditional view says that mathematics is as it appears: necessary and knowable
a priori. But the burden is to explain necessity and a priori knowledge clearly.

\subsection*{Whiteboard}

\[
\boxed{
\text{Traditional route: mathematics really is necessary and a priori.}
}
\]

\[
\boxed{
\text{Burden: explain necessity and a priority.}
}
\]

\section*{3. Kant's Attempted Answer}

Kant gives one of the most important historical attempts to explain the special status of
mathematics.

For Kant, arithmetic and geometry are \emph{synthetic a priori}. They are not merely matters
of definition, but they are still knowable independently of particular experience.

Kant's basic idea:

\[
\text{Mathematics concerns the forms through which we experience the world.}
\]

Geometry concerns the form of spatial intuition. Arithmetic concerns forms of spatial and
temporal intuition.

Mathematics is necessary because we cannot structure experience in any other way.
Mathematics is a priori because we do not need to inspect particular events in the world in
order to know these forms.

\subsection*{Teaching Point}

Kant tries to explain both sides of mathematics: its necessity and its connection to
experience. But his appeal to intuition creates further philosophical problems.

\subsection*{Whiteboard}

\[
\textbf{Kant:}
\]

\[
\text{Mathematics} = \text{synthetic a priori}
\]

\[
\text{Geometry} \rightarrow \text{space}
\]

\[
\text{Arithmetic} \rightarrow \text{space and time}
\]

\[
\text{Necessity comes from forms of intuition.}
\]

\section*{4. The Non-Traditional Route}

Some philosophers reject the traditional picture. They deny that mathematics is necessary
or a priori, or they sharply limit those notions.

This route is especially attractive to some empiricists and naturalists.

But this view has its own burden. If mathematics is not really necessary or a priori, then
why has it seemed that way to so many philosophers and mathematicians?

\subsection*{Teaching Point}

A philosopher may deny that mathematics is necessary or a priori. But then she must
explain why mathematics appears to have this special status.

\subsection*{Whiteboard}

\[
\boxed{
\text{Non-traditional route: mathematics is not necessary/a priori.}
}
\]

\[
\boxed{
\text{Burden: explain the appearance.}
}
\]

\section*{5. Global Matters: Objects and Objectivity}

The chapter then turns from knowledge to metaphysics and semantics.

A philosophy of mathematics must ask:

\[
\text{What is mathematics about?}
\]

Mathematical language appears to refer to objects:

\[
2,\quad 17,\quad \pi,\quad \mathbb{N},\quad \mathbb{R}
\]

The theorem that there are infinitely many primes seems to be about numbers.

So we ask:

\[
\boxed{
\text{Do mathematical objects exist?}
}
\]

\subsection*{Teaching Point}

Mathematics appears to have a subject matter. The philosopher must decide whether to
take that appearance at face value.

\subsection*{Whiteboard}

\[
\text{Mathematical discourse looks object-involving.}
\]

\[
\text{Numerals look like names.}
\]

\[
\text{What do they name?}
\]

\section*{6. Realism, Idealism, and Nominalism about Objects}

The chapter distinguishes three broad positions about mathematical objects.

\[
\begin{array}{c|c}
\textbf{View} & \textbf{Claim} \\
\hline
\text{Realism in ontology} &
\text{Mathematical objects exist objectively.} \\
\text{Idealism} &
\text{Mathematical objects depend on mind.} \\
\text{Nominalism} &
\text{Mathematical objects do not exist, or are linguistic constructions.}
\end{array}
\]

The realist says that at least some mathematical objects exist independently of
mathematicians, their minds, and their language.

The idealist says that mathematical objects depend on human mental activity.

The nominalist denies, in one way or another, that there are mathematical objects.

\subsection*{Teaching Point}

The question is not merely whether mathematicians talk as though numbers exist. The
question is whether such talk should be taken literally.

\subsection*{Whiteboard}

\[
\begin{array}{c|c}
\textbf{Realism} & \text{Numbers exist independently.} \\
\textbf{Idealism} & \text{Numbers depend on mind.} \\
\textbf{Nominalism} & \text{Numbers do not exist.}
\end{array}
\]

\section*{7. The Realist's Advantage and Problem}

Realism in ontology has an obvious advantage. If mathematical objects are eternal,
abstract, acausal, and independent of the physical world, then mathematical truths do not
depend on contingent facts.

This helps explain necessity.

\[
\text{Abstract objects}
\quad \Rightarrow \quad
\text{necessary truths?}
\]

But realism creates a deep epistemological problem.

If mathematical objects are outside space and time, and if they are causally inert, then how
can physical beings like us know anything about them?

\[
\text{Physical knower}
\quad \longrightarrow ? \quad
\text{abstract object}
\]

\subsection*{Teaching Point}

Realism gives a natural explanation of mathematical necessity, but it creates a mystery
about mathematical knowledge.

\subsection*{Whiteboard}

\[
\begin{array}{c|c}
\textbf{Realist advantage} & \textbf{Realist problem} \\
\hline
\text{explains necessity} & \text{explains knowledge?} \\
\text{objects independent} & \text{objects acausal} \\
\text{truth not contingent} & \text{how do we access them?}
\end{array}
\]

\section*{8. Truth-Value Realism}

The chapter distinguishes realism about objects from realism about truth.

\[
\textbf{Realism in ontology:}
\quad
\text{Mathematical objects exist objectively.}
\]

\[
\textbf{Realism in truth-value:}
\quad
\text{Mathematical statements have objective truth-values.}
\]

These two views often go together, but they are not the same.

A truth-value realist says that mathematical statements are objectively true or false,
independently of what mathematicians know, believe, or can prove.

\[
\text{Truth} \neq \text{knowability}
\]

\subsection*{Teaching Point}

One can ask two different questions: whether mathematical objects exist, and whether
mathematical statements are objectively true or false.

\subsection*{Whiteboard}

\[
\begin{array}{c|c}
\textbf{Ontology} & \textbf{Truth-value} \\
\hline
\text{Do numbers exist?} &
\text{Are mathematical statements objectively true or false?}
\end{array}
\]

\[
\text{Truth may outrun proof.}
\]

\section*{9. Anti-Realism in Truth-Value}

The anti-realist in truth-value denies that mathematical truth is wholly independent of the
mathematician.

There are two main forms:

\[
\begin{array}{c|c}
\textbf{Moderate anti-realism} &
\text{Truth depends on mind, proof, or construction.} \\
\textbf{Radical anti-realism} &
\text{Mathematical assertions lack genuine truth-values.}
\end{array}
\]

A moderate anti-realist may say:

\[
\text{If a mathematical statement is true, then it is knowable in principle.}
\]

This rejects the realist separation between truth and knowability.

Some anti-realists also reject bivalence:

\[
P \lor \neg P
\]

That is, they deny that every unambiguous mathematical statement must be either
determinately true or determinately false.

\subsection*{Teaching Point}

Truth-value anti-realism connects metaphysics, semantics, epistemology, and logic. If
truth depends on knowability, then classical logic may come under pressure.

\subsection*{Whiteboard}

\[
\textbf{Truth-value realism:}
\quad
\text{Truth can outrun knowledge.}
\]

\[
\textbf{Truth-value anti-realism:}
\quad
\text{Truth is tied to knowability.}
\]

\[
\text{Bivalence?}
\quad
P \lor \neg P
\]

\section*{10. Benacerraf's Dilemma}

The chapter introduces the central dilemma from Paul Benacerraf.

We want mathematical language to be continuous with ordinary and scientific language.
Numerals look like names. Mathematical sentences look like ordinary truth-bearing
sentences.

This pushes us toward realism:

\[
\text{Uniform semantics}
\quad \Rightarrow \quad
\text{mathematical objects and objective truth}
\]

But realism creates an epistemological problem:

\[
\text{Abstract, acausal objects}
\quad \Rightarrow \quad
\text{How can we know them?}
\]

So the dilemma is:

\[
\boxed{
\text{Good semantics seems to produce bad epistemology.}
}
\]

\[
\boxed{
\text{Good epistemology seems to require non-standard semantics.}
}
\]

\subsection*{Teaching Point}

Benacerraf's dilemma is one of the central organizing problems in contemporary
philosophy of mathematics.

\subsection*{Whiteboard}

\[
\begin{array}{c|c}
\textbf{Semantic desideratum} & \textbf{Epistemological problem} \\
\hline
\text{Math language works like ordinary language} &
\text{How do we know abstract objects?} \\
\text{Numerals refer to numbers} &
\text{Numbers are causally inert} \\
\text{Mathematical statements are true} &
\text{No causal contact}
\end{array}
\]

\[
\boxed{\text{Benacerraf: semantics vs. epistemology}}
\]

\section*{11. Mixed Views}

The chapter notes that realism about ontology and realism about truth-value can come
apart.

The four possible combinations are:

\[
\begin{array}{c|c|c}
& \textbf{Truth-value realism} & \textbf{Truth-value anti-realism} \\
\hline
\textbf{Ontological realism} &
\text{Objects exist; truths objective} &
\text{Objects exist; truth tied to knowability} \\
\hline
\textbf{Ontological anti-realism} &
\text{No objects; truths objective} &
\text{No objects; truth not objective}
\end{array}
\]

The most natural alliances are:

\[
\text{ontology realism} + \text{truth-value realism}
\]

and

\[
\text{ontology anti-realism} + \text{truth-value anti-realism}.
\]

But mixed views are possible. For example, one might deny mathematical objects while
still trying to preserve objective mathematical truth.

\subsection*{Teaching Point}

Do not collapse the question of objects into the question of truth. A philosopher may be
realist about one and anti-realist about the other.

\subsection*{Whiteboard}

\[
\text{Objects?}
\quad \neq \quad
\text{Objective truth?}
\]

\section*{12. The Mathematical and the Physical}

The chapter next turns to the application of mathematics.

Mathematics is deeply involved in science. This is not limited to what is officially called
``applied mathematics.'' Geometry, probability, group theory, logic, and analysis all
illuminate parts of nature.

The central problem is:

\[
\boxed{
\text{How can mathematics explain physical phenomena?}
}
\]

A physical event may be explained by reference to a mathematical fact. For example, a
pattern on an oscilloscope might be explained by saying that a function solving a
differential equation has a zero at a certain value.

But why should that count as an explanation?

\subsection*{Teaching Point}

It is not enough to observe that mathematics is useful in science. A philosophy of
mathematics should explain how mathematics can be useful in science.

\subsection*{Whiteboard}

\[
\text{Zero of a function}
\quad \longrightarrow ? \quad
\text{pattern on oscilloscope}
\]

\[
\text{Mathematical fact}
\quad \longrightarrow ? \quad
\text{physical explanation}
\]

\section*{13. Three Problems of Applicability}

The chapter distinguishes several problems about the application of mathematics.

\[
\begin{array}{c|p{0.62\textwidth}}
\textbf{Semantic problem} &
How do mixed mathematical/physical statements have meaning? \\

\textbf{Metaphysical problem} &
How do mathematical objects, if there are any, relate to physical objects? \\

\textbf{Conceptual problem} &
Why are specific mathematical concepts and structures useful in describing reality?
\end{array}
\]

Examples:

\[
\text{Why is arithmetic useful?}
\]

\[
\text{Why is group theory useful for symmetry?}
\]

\[
\text{Why are Hilbert spaces useful in physics?}
\]

\subsection*{Teaching Point}

The problem of applicability is not one problem but a cluster of problems. We need to
explain language, ontology, modeling, and scientific practice.

\subsection*{Whiteboard}

\[
\text{Applicability}
=
\begin{cases}
\text{semantic} \\
\text{metaphysical} \\
\text{conceptual}
\end{cases}
\]

\section*{14. Indispensability and Its Limits}

A popular argument for realism appeals to mathematics' role in science.

\[
\text{Mathematics is indispensable to science.}
\]

\[
\text{Science is approximately true.}
\]

\[
\therefore \text{Mathematics is true.}
\]

This is the basic shape of the indispensability argument.

But the chapter warns that this is too quick. Even if mathematics is indispensable, we still
need to explain how mathematics is applied.

\[
\text{Indispensability itself needs explanation.}
\]

\subsection*{Teaching Point}

The indispensability argument cannot simply point to the success of science and stop. It
must explain how mathematical statements contribute to scientific knowledge.

\subsection*{Whiteboard}

\[
\begin{array}{c}
\text{Math indispensable to science} \\
\text{Science true or approximately true} \\
\hline
\text{Math true?}
\end{array}
\]

\[
\text{Question: How does the application work?}
\]

\section*{15. The Unreasonable Effectiveness Problem}

The chapter also notes a deeper puzzle: pure mathematics often finds unexpected
application long after it is developed.

Mathematicians develop structures for internal mathematical reasons. Later, physicists or
scientists discover that these structures fit empirical reality.

This creates the famous sense that mathematics is ``unreasonably effective.''

\subsection*{Teaching Point}

The problem is not merely that mathematics applies. The problem is that mathematics
often applies in unexpected ways.

\subsection*{Whiteboard}

\[
\text{Pure mathematics}
\quad \longrightarrow
\quad \text{later scientific application}
\]

\[
\text{Why was the mathematics already there?}
\]

\section*{16. Local Matters: Theorems, Theories, and Concepts}

The final section shifts from global questions to local philosophical problems.

Not all philosophy of mathematics asks broad questions like:

\[
\text{Do numbers exist?}
\]

Some philosophical questions arise from particular mathematical results, theories, and
concepts.

\[
\text{Local results can reopen global questions.}
\]

Examples include:

\[
\text{Skolem paradox}
\]

\[
\text{independence results in set theory}
\]

\[
\text{G\"odel's incompleteness theorem}
\]

\[
\text{the development of the concept of function}
\]

\subsection*{Teaching Point}

Philosophy of mathematics is not only about large-scale metaphysical theories. It also
interprets specific mathematical concepts and results.

\subsection*{Whiteboard}

\[
\text{Local mathematics}
\quad \longrightarrow
\quad \text{global philosophy}
\]

\section*{17. The Skolem Paradox}

The Skolem paradox arises from results in first-order model theory.

Roughly: if a first-order theory has an infinite model, then it has models of many different
infinite sizes.

This creates a puzzle. First-order set theory can have countable models, even though set
theory proves that there are uncountably many real numbers.

This is not a contradiction, but it raises philosophical questions:

\[
\text{Do our formal theories determine their intended subject matter?}
\]

\[
\text{Can first-order language capture the natural numbers or the real numbers uniquely?}
\]

\[
\text{How do we communicate determinate mathematical concepts?}
\]

\subsection*{Teaching Point}

The Skolem paradox raises questions about reference, determinacy, and whether formal
theories fully capture mathematical concepts.

\subsection*{Whiteboard}

\[
\textbf{Skolem paradox}
\]

\[
\text{First-order theories have unintended models.}
\]

\[
\therefore \quad
\text{Can formal language fix reference?}
\]

\section*{18. Independence and the Continuum Hypothesis}

Zermelo--Fraenkel set theory with choice, or ZFC, is a powerful standard theory of sets.

But some important questions cannot be decided by ZFC. The most famous example is the
continuum hypothesis.

\[
\textbf{Continuum Hypothesis:}
\]

\[
\text{There is no infinite cardinality strictly between } |\mathbb{N}| \text{ and } |\mathbb{R}|.
\]

Neither the continuum hypothesis nor its negation can be proved from the standard
axioms of ZFC.

This raises philosophical questions:

\[
\text{Does CH have a determinate truth-value?}
\]

\[
\text{Should we search for new axioms?}
\]

\[
\text{Or may we choose whichever extension is most useful?}
\]

\subsection*{Teaching Point}

Independence results put pressure on truth-value realism. If a statement is undecidable
from accepted axioms, we must ask whether it is still objectively true or false.

\subsection*{Whiteboard}

\[
\textbf{CH independent of ZFC}
\]

\[
\text{Realist: CH is true or false; find new axioms.}
\]

\[
\text{Anti-realist: maybe no fact of the matter.}
\]

\section*{19. G\"odel's Incompleteness Theorem}

G\"odel's incompleteness theorem concerns formal theories of arithmetic.

Roughly:

\[
\text{Any sufficiently rich, consistent, effective theory of arithmetic is incomplete.}
\]

That is, there will be a sentence \(P\) such that neither \(P\) nor \(\neg P\) is derivable in
the theory.

\[
T \nvdash P
\qquad
T \nvdash \neg P
\]

A truth-value realist may say:

\[
\text{Truth outruns proof in any fixed formal system.}
\]

A truth-value anti-realist may say:

\[
\text{Perhaps undecidable statements lack determinate truth-values.}
\]

\subsection*{Teaching Point}

Incompleteness forces a distinction between formal derivability and mathematical truth.

\subsection*{Whiteboard}

\[
\textbf{G\"odel}
\]

\[
\text{Truth} \neq \text{provability in a fixed formal system}
\]

\[
T \nvdash P
\qquad
T \nvdash \neg P
\]

\section*{20. Concepts: Function, Polyhedron, and \(dx\)}

The chapter ends by emphasizing conceptual interpretation.

Mathematical development often changes or clarifies what a concept means.

Examples:

\[
\text{function}
\]

\[
\text{polyhedron}
\]

\[
dx
\]

The modern concept of function emerged through foundational work in analysis.
Lakatos's work on polyhedra shows how proofs and refutations can reshape a concept.
The history of calculus shows that expressions like \(dx\) and \(dy/dx\) required long
interpretive clarification.

\subsection*{Teaching Point}

Sometimes mathematics can proceed before its concepts are fully interpreted. But tensions
within mathematics often force philosophical clarification.

\subsection*{Whiteboard}

\[
\text{Mathematics often asks:}
\]

\[
\text{What does this concept really mean?}
\]

\[
\text{function} \quad \text{polyhedron} \quad dx
\]

\section*{Conclusion: The Chapter's Map of Philosophy of Mathematics}

The chapter gives us a map of the major questions any philosophy of mathematics must
address.

\[
\begin{array}{c|p{0.62\textwidth}}
\textbf{Issue} & \textbf{Question} \\
\hline
\text{Necessity} &
Why do mathematical truths seem necessary? \\

\text{A priority} &
How can mathematical knowledge be independent of experience? \\

\text{Objects} &
Do numbers, sets, functions, and points exist? \\

\text{Truth} &
Are mathematical statements objectively true or false? \\

\text{Knowledge} &
How can we know mathematical truths? \\

\text{Application} &
Why does mathematics apply to the physical world? \\

\text{Local interpretation} &
What do particular theorems, theories, and concepts show?
\end{array}
\]

\section*{Final Framing}

To teach this chapter well, present it as a map of the problems rather than as a list of
answers.

A philosophy of mathematics must explain why mathematics seems necessary and
a priori, what mathematical language is about, whether mathematical truth is objective,
how mathematical knowledge is possible, why mathematics applies to science, and how
specific mathematical results reshape philosophical questions.

\[
\boxed{
\text{The chapter gives the problem-space of philosophy of mathematics.}
}
\]

\[
\boxed{
\text{Later chapters explore competing solutions.}
}
\]

\section*{Discussion Questions}

\begin{enumerate}[leftmargin=2em]
    \item Is mathematical knowledge more certain than scientific knowledge?
    \item Could \(7+5=12\) have been false?
    \item Is mathematics really a priori, or does it only appear that way?
    \item Do numbers exist independently of human minds?
    \item Can mathematical statements be objectively true even if mathematical objects do not exist?
    \item Does truth outrun proof in mathematics?
    \item How can an abstract mathematical fact explain a physical event?
    \item Does the success of mathematics in science support realism?
    \item Does the continuum hypothesis have a determinate truth-value?
    \item What does G\"odel's incompleteness theorem show about the relation between truth and proof?
\end{enumerate}

\section*{Key Terms}

\begin{description}[leftmargin=2em]
    \item[Necessity] The modal status of something that could not have been otherwise.

    \item[A priori knowledge] Knowledge independent of experience with the specific course of events in the actual world.

    \item[A posteriori knowledge] Knowledge based on experience.

    \item[Realism in ontology] The view that at least some mathematical objects exist objectively and independently of mathematicians.

    \item[Idealism] The view that mathematical objects depend on the mind.

    \item[Nominalism] The view that mathematical objects do not exist, or that apparent reference to them must be reinterpreted.

    \item[Realism in truth-value] The view that mathematical statements have objective truth-values independent of minds, language, or convention.

    \item[Anti-realism in truth-value] The view that mathematical truth depends in some way on minds, proof, construction, or knowability.

    \item[Bivalence] The principle that every unambiguous statement is either true or false.

    \item[Benacerraf's dilemma] The problem that a semantics continuous with ordinary language seems to require mathematical objects, while such objects create a deep epistemological mystery.

    \item[Applicability problem] The problem of explaining how mathematics applies to and explains the physical world.

    \item[Indispensability argument] The argument that mathematics should be accepted as true because it is indispensable to our best science.

    \item[Skolem paradox] The puzzle that first-order theories can have unintended models, including models of unexpected cardinality.

    \item[Continuum hypothesis] The claim that there is no infinite cardinality strictly between the size of the natural numbers and the size of the real numbers.

    \item[Incompleteness theorem] G\"odel's result that any sufficiently rich, consistent, effective formal theory of arithmetic leaves some arithmetic sentence undecided.
\end{description}

\end{document}